% page 244 of the facsimile (image p-243 of the scan)
% block 162.6 x 254.9 mm ; left margin 8.0 mm
\pgc{-12.700mm}{0.000mm}{
\l{}
\l{}
\l{}
\l{\cel{2}236}
\l{}
\l{inkandecar. (netrans.)\cel{2}Qua esas tante varmigita ke lu}
\l{\cel{3}aspektas blanka. - EFIS.}
\l{}
\l{inkardinar. (trans.)\cel{2}Destinar, perpetue ed absolute, ula}
\l{\cel{3}kleriko a la servo di diocezo determinata.}
\l{}
\l{inkarnacar.\cel{2}(trans.)\cel{2}(\sou{teol.}) Glitar su aden korpo ek karno}
\l{\cel{3}homatra. - DEFIS.}
\l{}
\l{\sou{inkasar}. (\sou{trans}.)Recevar e pozar aden la kaso (di la afero)}
\l{\cel{3}pekunio-quanto. - DFI.}
\l{}
\l{inkastrar (trans.)\cel{2}Inserar (objekto) aden altra objekto qua}
\l{\cel{3}entaliesis por recevar olu ; fixigar (lapido) aden la kas-}
\l{\cel{2}tono, per abatar la rebordo di olca adcirkum la lapido.-FIS.}
\l{}
\l{inklinar. (trans., ad). Igar lua axo kelke deviacar de lua}
\l{\cel{3}vertikalo. ( EFIS.}
\l{}
\l{inkluzar. (trans., (ad)en.)\cel{2}Enpozar,cirkumklozar. - DEFIRS.}
\l{}
\l{inkoativa.\cel{2}(gram.)\cel{2}Qua expresas la komenco di ago.-DEFIS.}
\l{}
\l{inkognito.\cel{2}Sen esar konocata,remarkata.- DEFIRS.}
\l{}
\l{\sou{inkombrar}. (\sou{trans., ye)}.\cel{2}(\sou{aludante amaso de kozi)}.\cel{2}\sou{Embarasa}r}
\l{\cel{3}per esar obstaklo pri la cirkulo. - EFIS.}
\l{}
\l{inkrustar. (trans.)\cel{2}Kovrar (objekto) ye strato petra e krus-}
\l{\cel{3}tatra. - DEFIRS.}
\l{}
\l{inkubo.\cel{2}(teolo.)\cel{2}Demono qua kopulacis, lokizinte su sur la}
\l{\cel{3}dormanto. - DEFIS.}
\l{}
\l{\sou{inkubacar}. (\sou{trans.)}\cel{3}Havar en su la jermo di morbo, qua kres-}
\l{\cel{3}kas til kande onu maladeskas. - DEFIS.}
\l{}
\l{inkulkar. (trans.) Igar par-enirar en la spirito. - EFIS.}
\l{}
\l{inkunablo.\cel{2}Libro qua evas de la epoko kande onu jus inventabis}
\l{\cel{3}la imprimarto (generale til la yaro 1500). - DEFIS.}
\l{}
\l{\sou{inkurar}. (\sou{trans.})\cel{2}Igar su en la kazo subisor. F.}
\l{}
\l{inkursar. (netrans.) I. Devast-irupto da militanti aden la lando}
\l{\cel{3}enemika. - II. (\sou{metaf.})\cel{2}Trakto acesora di temo qua ne}
\l{\cel{3}apartenas a la domeno precipua. - EFIRS.}
\l{}
\l{inocenta. (pri). Qua ne agis ulo mala. - EFIS.}
\l{}
\l{\sou{inokular}. (\sou{trans.}) (\sou{medic.})\cel{2}Transdonar ad\cel{1}\sou{ulu}\cel{1}(viruso, jermodi}
\l{\cel{3}morbo kontagiala) ; o transdonar artificale la viruso (fe-}
\l{\cel{3}bligita) di morbo, kom prezervivo kontre la ipsa morbo.}
\l{\cel{3}- DEFIRS.}
}
